\documentclass[10pt,a4paper]{amsart}
\usepackage[margin=19mm]{geometry}
\usepackage{amsmath,amssymb,mathtools}
\usepackage[scr=rsfs,cal=euler]{mathalfa}
\usepackage{xfrac}
\usepackage[unicode,hidelinks,bookmarks=false]{hyperref}
\newcommand{\ie}{i.e.\ }
\newcommand{\cf}{cf.\ }
\newcommand{\resp}{respectively\ }
\newcommand{\Subsection}{\subsection}
\providecommand{\nobreakdash}{\nobreak\hbox{-}}
\setlength{\emergencystretch}{2em}
\multlinegap0pt
\usepackage{amssymb}
\usepackage{csquotes}
\newtheorem{theorem}{Theorem}
\title{On the Role of the $Ky~Fan$ Metric in Rough Ideal Convergence in Probability}
\author{Tamim Aziz, Sanjoy Ghosal}
\date{Source: arXiv:2601.01797v1 (CC BY 4.0)}
\begin{document}
\maketitle

\noindent\emph{Source and licence.} On the Role of the $Ky~Fan$ Metric in Rough Ideal Convergence in Probability. Version arXiv:2601.01797v1 (CC BY 4.0), distributed by arXiv under the Creative Commons Attribution 4.0 International licence, http://creativecommons.org/licenses/by/4.0/, as recorded in the arXiv licence metadata for this exact version and shown on the abstract page https://arxiv.org/abs/2601.01797v1. Full licence notice: \texttt{licenses/arxiv-2601.01797-CC-BY-4.0.txt}.\par\smallskip

\noindent\textit{Editorial scope.} Counted unit: the article's theorem th (source0.tex lines 646--648). The article's complete original proof of it is delivered with it (source0.tex lines 649--697, one \texttt{proof} environment of 49 lines). No mathematics is written, repaired or re-derived: the statement and the proof are the article's own lines, in the article's own notation, and no retained line is substituted or re-flowed.  No further source-local context is needed: every cross-reference of every retained line resolves inside the two delivered spans.  Nothing was omitted from any retained span, so the package carries no omission record.  No retained line contains a citation, so the wrapper carries no bibliography and no bibliography entry.  No retained line contains a figure environment, an image-inclusion command or drawing code, so no image is delivered and the package declares no asset dependency.  Editorial formatting only, in addition: an AMS \texttt{amsart} preamble that restates the article's own declarations and macros used by the retained lines, namely the environments \texttt{theorem} on separate counters, together with the standard packages those lines need.  The retained environments print in this document as Theorem 1 in order of appearance, whereas the article prints the same items with the numbering of its own full document.  The complete native source of the article as distributed in the arXiv e-print for this exact version is bundled unchanged as \texttt{sources/192112/source0.tex}.
\begin{theorem}\label{maximal th}
Suppose $(\mathfrak{X}^0, \rho)$ has at least two distinct elements, and $\mathcal{I}$ is an admissible ideal. Then $\mathcal{I}$ is maximal if and only if $\mathcal{I}^{\mathbb{P}}$-$LIM^rX_i=\Gamma^{r^s}_{\underline{X}}(\mathcal{I}^{\mathbb{P}})$ holds for each $r\geq 0$ and each $\underline{X}$ in $\mathfrak{X}$.
\end{theorem}

\begin{proof}
	Assume that $\mathcal{I}$ is a maximal admissible ideal on $\mathbb{N}$. Let $\varepsilon,\delta>0$ be arbitrary.
Pick any $r\geq 0$ and any sequence $\underline{X}$ with values in $\mathfrak{X}$. Then, for each $X_*\in\mathcal{I}^{\mathbb{P}}$-$LIM^rX_i$, we have 
\begin{align*}
	&\{n\in\mathbb{N}:\mathbb{P}(d(X_n,X_*)\geq r+\varepsilon)>\delta\}\in\mathcal{I}\\
	\Rightarrow~& \{n\in\mathbb{N}:\mathbb{P}(d(X_n,X_*)\geq r+\varepsilon)\leq\delta\}\notin\mathcal{I}~\mbox{ (since $\mathcal{I}$ is non-trivial)}\\
	\Rightarrow~& \{n\in\mathbb{N}:1-\mathbb{P}(d(X_n,X_*)< r+\varepsilon)\leq\delta\}\notin\mathcal{I}\\
\Rightarrow~& \{n\in\mathbb{N}:\mathbb{P}(d(X_n,X_*)< r+\varepsilon)\geq 1-\delta\}\notin\mathcal{I}.
\end{align*}
This ensures that $X_*\in \Gamma^{r^s}_{\underline{X}}(\mathcal{I}^{\mathbb{P}})$, i.e., $\mathcal{I}^{\mathbb{P}}$-$LIM^rX_i\subseteq\Gamma^{r^s}_{\underline{X}}(\mathcal{I}^{\mathbb{P}})$.\\

On the other hand, pick any $Y\in\Gamma^{r^s}_{\underline{X}}(\mathcal{I}^{\mathbb{P}})$. Then observe that 
\begin{align*}
	&\{n\in\mathbb{N}:\mathbb{P}(d(X_n,X_*)< r+\varepsilon)\geq 1-\delta\}\notin\mathcal{I}\\
	\Rightarrow~& \{n\in\mathbb{N}:\mathbb{P}(d(X_n,X_*)< r+\varepsilon)<1-\delta\}\in\mathcal{I}~\mbox{ (since $\mathcal{I}$ is admissible and maximal)}\\
	\Rightarrow~& \{n\in\mathbb{N}:1-\mathbb{P}(d(X_n,X_*)\geq r+\varepsilon)<1-\delta\}\in\mathcal{I}\\
	\Rightarrow~& \{n\in\mathbb{N}:\mathbb{P}(d(X_n,X_*)\geq r+\varepsilon)>\delta\}\in\mathcal{I}.
\end{align*}
Consequently, we obtain $Y\in\mathcal{I}^{\mathbb{P}}$-$LIM^rX_i$, i.e., $\Gamma^{r^s}_{\underline{X}}(\mathcal{I}^{\mathbb{P}})\subseteq \mathcal{I}^{\mathbb{P}}$-$LIM^rX_i$. Hence we can deduce that $\mathcal{I}^{\mathbb{P}}$-$LIM^rX_i=\Gamma^{r^s}_{\underline{X}}(\mathcal{I}^{\mathbb{P}})$.\\

	Conversely, suppose that $\mathcal{I}^{\mathbb{P}}$-$LIM^rX_i=\Gamma^{r^s}_{\underline{X}}(\mathcal{I}^{\mathbb{P}})$ holds for each $r\geq 0$ and each $\underline{X}$ in $\mathfrak{X}$.
Assume, on the contrary, that $\mathcal{I}$ is not maximal. Then there exists $A\subset\mathbb{N}$ such that either $A,\mathbb{N}\setminus A\in\mathcal{I}$ or $A,\mathbb{N}\setminus A\notin\mathcal{I}$. Since $\mathcal{I}$ is non-trivial, we must have $A,\mathbb{N}\setminus A\notin\mathcal{I}$. By the given hypothesis, there exist $Y,Z\in \mathfrak{X}^0$ such that $\alpha=\rho(Y,Z)>0$. Since the infimum in the definition of $\rho(Y,Z)=\inf \{\varepsilon>0:\mathbb{P}(d(Y,Z)>\varepsilon)\leq\varepsilon\}$ is attained, we have 
$$
\mathbb{P}(d(Y,Z)>\alpha)\leq\alpha.
$$
We fix any $0<\beta<\alpha$. 
Now, pick any $\varepsilon>0$ such that $\beta+\varepsilon<\alpha$.
Then, we will have 
$$
\mathbb{P}(d(Y,Z)>\beta+\varepsilon)>\beta+\varepsilon~\mbox{ (using the minimality of $\alpha$)}.
$$
Let us define $\underline{X}$ in $\mathfrak{X}$ as follows:
\begin{align*}
	X_n=
	\begin{cases}
		Y~&\mbox{ if }~ n\in A,\\
		Z ~&\mbox{ if }~n\in \mathbb{N}\setminus A.
	\end{cases}
\end{align*}
Then, observe that 
$$
A\subseteq\{n\in\mathbb{N}:\mathbb{P}(d(X_n,Y)<\beta+\varepsilon)>1-\delta\}\notin\mathcal{I}~\mbox{ for every }~\delta>0,
$$
i.e., $Y\in\Gamma^{\beta^s}_{\underline{X}}(\mathcal{I}^{\mathbb{P}})$. Also, observe that 
$$
\mathbb{N}\setminus A\subseteq\{n\in\mathbb{N}:\mathbb{P}(d(X_n,Y)>\beta+\varepsilon)>\beta+\varepsilon\}\notin\mathcal{I},
$$
i.e., $Y\notin \mathcal{I}^{\mathbb{P}}$-$LIM^{\beta}X_i$ - which is a contradiction. Hence we deduce that $\mathcal{I}$ is maximal.
\end{proof}

\end{document}
